\documentclass[11pt]{article}
\usepackage[a4paper,margin=25mm]{geometry}
\usepackage[T1]{fontenc}
\usepackage{lmodern,microtype,amsmath,amssymb,amsthm,booktabs,array,tabularx}
\usepackage[round,authoryear]{natbib}
\usepackage[colorlinks=true,linkcolor=blue,citecolor=blue,urlcolor=blue]{hyperref}
\setlength{\parindent}{1.2em}
\setlength{\parskip}{3pt}
\setlength{\emergencystretch}{2em}
\newtheorem{theorem}{Theorem}
\newtheorem{proposition}[theorem]{Proposition}
\newcolumntype{Y}{>{\raggedright\arraybackslash}X}
\newcolumntype{L}[1]{>{\raggedright\arraybackslash}p{#1}}
\newcommand{\E}{\mathbb{E}}
\newcommand{\Prob}{\mathbb{P}}
\newcommand{\ind}{\mathbf{1}}
\hypersetup{pdftitle={Cortex: Sleep-Inspired Memory Consolidation for Language Agents},pdfauthor={Brendan Playford}}
\title{Cortex: Sleep-Inspired Memory Consolidation\\for Language Agents}
\author{Brendan Playford}
\date{October 2026}
\begin{document}
\maketitle
\begin{abstract}
Sleep-dependent memory research and complementary learning systems theory motivate a distinction between rapid episodic capture and slower integration. We use this distinction to design Cortex, a memory engine for language agents that preserves source observations while consolidating derived memories. The analogy is functional: current consolidation uses explicit model calls and validated database transitions; it neither simulates sleep stages nor implements generative dreaming. Retrieval combines lexical and semantic evidence, and an optional reader checks literal quotations and computes source-linked date differences. Conditional proofs describe scope preservation, atomic acknowledgement, activation decay, and source coverage in an abstract state machine. An exploratory study uses 2,814 private user messages and sixty held-out public memory cases. Source-assisted answering matches 47 references, including eleven conditional answers, versus 45 for a direct baseline; the difference does not establish an advantage. Locally trained retention and memory-type classifiers build on a pretrained sentence encoder. Their best accuracies against provisional teacher labels are 64.2\% and 91.9\%, respectively, below constant-label accuracy checks on the imbalanced test set. The joint rule accepts 12.4\% of candidates with 28.9\% disagreement, leaving classification advisory. We distinguish these measured results from vendor-reported platform scores and specify a comparative protocol with frozen datasets, matched readers and budgets, human-reviewed labels, and separate tests of consolidated memory, source retrieval, and replay. The protocol makes the sleep analogy testable without treating synthetic recollections as observed events.
\end{abstract}

\section{Introduction}
Language agents frequently need information absent from their current context: a user's preference, a correction to an earlier statement, or a detail from a previous task. External memory can retain such information without changing model parameters. Its reliability depends on separate decisions: what to retain, which source supports it, whether a later statement supersedes it, and which evidence to retrieve for a new query.

Summary-based memory introduces an information bottleneck. A fluent summary may omit a detail that later becomes relevant. Improving retrieval over those summaries cannot recover that detail. Conversely, retrieving original conversations exposes the reader to quotations, obsolete statements, and irrelevant requests. A useful memory system needs a compact representation and access to attributed evidence, with explicit lifecycle rules for each.

This work gives a common state model for observations, derived memories, and source excerpts; separates statistical decisions from validated persistence; and implements a private evaluation and training workflow. The empirical study is exploratory. Its purpose is to identify failure mechanisms and establish an auditable baseline. We also separate locally trained classification parameters from external model capabilities, reanalyse class imbalance, and specify the evidence needed for a credible comparison with established memory platforms.

\section{Biological motivation and scope}
Complementary learning systems theory explains how rapid episodic learning and slower integration can coexist without continually overwriting structured knowledge \citep{cls}. This motivates separating immutable observations from derived memories. It does not imply that database rows instantiate hippocampal circuits or that summarization is equivalent to synaptic consolidation.

Sleep research supplies more specific hypotheses. Human intracranial recordings during targeted memory reactivation associate successful NREM consolidation with coordinated hippocampal and cortical activity \citep{nrem2025}. In mice, selective disruption of small adult-born neuronal ensembles during REM sleep provides evidence for a causal contribution to associative fear-memory consolidation \citep{rem2025}. These results concern distinct experimental settings and do not establish a single universal REM consolidation rule.

A computational model by \citet{dreaming} uses complementary waking, perturbed NREM replay, and generative REM phases for visual representation learning. Its dream phase generates novel combinations of prior experience. That study evaluates visual representations; transfer of the mechanism to conversational memory remains an empirical question. Cortex currently implements attributed capture, explicit offline integration, retrieval, activation decay and optional classifier training. It does not implement a biological sleep cycle or a generative replay phase. Its formal results below concern software invariants, not a theorem derived from neuroscience.

Table~\ref{tab:sleep-mapping} separates the implemented mechanisms from hypotheses. The word \emph{sleep-inspired} names the architectural motivation; it does not claim biological fidelity. Model inference over a database is not neuronal replay, and relevance decay is not a measurement of synaptic weakening.

\begin{table}[htbp]
\centering\small
\begin{tabularx}{\linewidth}{@{}L{.22\linewidth}Y L{.19\linewidth}@{}}
\toprule
Biological motivation & Software interpretation & Status\\
\midrule
Rapid episodic capture & Store attributed observations before deriving reusable memories & Implemented\\
Slower integration & Consolidate observations with explicit source support and contradiction history & Implemented\\
NREM replay & Revisit source-grounded training examples under a fixed replay budget & Proposed experiment\\
REM recombination & Generate recombined training examples, separate from observed history & Proposed experiment\\
Selective persistence & Reduce retrieval activation with time and preserve reinforced memories & Engineering heuristic\\
\bottomrule
\end{tabularx}
\caption{Functional correspondences and their limits. No row establishes a neural mechanism or a simulated sleep stage.}
\label{tab:sleep-mapping}
\end{table}

The analogy suggests a testable next experiment: compare direct classification training, replay of source-grounded examples, and recombined synthetic replay at matched data and compute budgets. Source attribution and correction tests must accompany transfer measurements, because a plausible synthetic recollection is not evidence that an event occurred. Generated examples would remain a separate training artifact and could never become autobiographical observations without independent source support.

\section{Related work}
Cognitive Architectures for Language Agents describes modular memory and decision procedures within an agent architecture \citep{coala}. MemGPT uses hierarchical memory and control flow to manage a context window \citep{memgpt}. These works motivate treating memory as more than a collection of vector embeddings. Cortex concentrates on persistence and evidence. Scheduling, user authorization, and the interaction loop belong to its host.

LongMemEval separates indexing, retrieval, and reading, and evaluates information extraction, multi-session reasoning, temporal reasoning, updates, and abstention \citep{longmem}. This separation guides our experiments: answer failure alone does not identify the responsible component. Our small public probe is not comparable to a complete benchmark evaluation.

Sentence-BERT supplies a practical representation for semantic similarity \citep{sbert}. SetFit adapts a sentence encoder through contrastive learning before fitting a classification head \citep{setfit}. We use these methods as baselines for classifying source spans. Temperature scaling is selected on a development partition to examine probability calibration \citep{calibration}.

Memory-R1 learns memory operations using downstream rewards \citep{memoryr1}. MemRL adapts episodic-memory utility without updating the reasoning model's weights \citep{memrl}. Both directions make feedback quality consequential. We keep outcome utility, supervised labels, and model weights distinct. Grading discrepancies in our public probe provide a reason to audit rewards before optimizing a memory policy against them.

\section{Problem formulation}
An observation is a tuple
\[
 o=(i,u,x,t_o,\rho),
\]
where $i$ is its identifier, $u$ its owner scope, $x$ immutable source text, $t_o$ observation time, and $\rho$ provenance such as speaker and conversation identity. Observation time does not establish the date of an event described in $x$.

A derived memory is
\[
 m=(j,u,s,k,c,S,t_\ell,n,z),
\]
where $s$ is a summary, $k$ a semantic type, $c\in[0,1]$ evidence confidence, $S$ a nonempty set of supporting observation identifiers, $t_\ell$ the last supporting observation time, $n$ a reinforcement count, and $z$ lifecycle metadata. The supported types are fact, behavior, preference, relationship, project, and routine. They describe an annotation policy; they are not mutually exclusive properties of natural language in general.

An atomic candidate is an exact span $x[a:b]$ with half-open Unicode code-point offsets. Atomicity is an annotation goal, not a property guaranteed by span validation. Retention belongs to
\[
 \mathcal{R}=\{\text{durable},\text{temporary},\text{ignore},\text{uncertain}\}.
\]
Retention and type answer different questions. A one-time project request can be temporary. A quotation can resemble a preference without supporting an inference about the user's preference.

For query $q$ and evidence budget $B$, a conceptual objective is
\[
 \max_C \E[\mathrm{answer\ quality}\mid q,C]
 \quad\text{subject to}\quad
 \mathrm{tokens}(C)\leq B,\quad C\subseteq\mathcal{E}_u,
\]
where $\mathcal{E}_u$ is currently eligible evidence for owner $u$. Rank fusion approximates selection; scope and lifecycle constraints are enforced separately. Quality must distinguish unsupported retention, omitted facts, stale answers, attribution mistakes, and reader errors.

\section{Memory operations and retrieval}
\subsection{Consolidation}
Consolidation processes a bounded snapshot of pending observations. A model proposes insertion, reinforcement, or contradiction operations. Validation requires known operation types, bounded summaries, finite confidence values, sources within the current chunk, and update targets in the same owner scope. Contradiction records a successor while preserving the earlier memory's provenance.

Accepted operations and acknowledgement of their exact source snapshot commit in one transaction. Inference occurs outside that transaction. A malformed successful response may receive one format-repair request; the result passes the same validation. This bounds recovery without treating provider calls as exactly-once events.

Inference receives short request-local identifiers for observations and existing memories. An injective mapping $\pi$ associates each permitted stored identifier with a local alias; the application resolves $\pi^{-1}$ before validating or committing a proposal. Aliases are chosen outside the set of real identifiers to avoid collisions, and duplicate sources after resolution are rejected. Historical source lists are retained in storage and omitted from the model context, since they cannot support a new operation. These measures reduce identifier-copying burden without expanding the set of permitted sources.

The host decides when consolidation runs and whether observations may be sent to an embedding provider or classifier. Ordinary retrieval does not create a training label.

\subsection{Summary and source retrieval}
A query can retrieve derived memories or source excerpts. Lexical overlap and cosine similarity produce ranked lists within each representation. Reciprocal rank fusion assigns
\[
 R(d)=\sum_{c\in\mathcal{C}}
 \frac{\ind\{d\text{ appears in }c\}}{\kappa+\operatorname{rank}_c(d)},
 \qquad \kappa=60,
\]
following \citet{rrf}. The score orders candidates; it is not a relevance probability.

The default segmentation divides source text into overlapping windows of $W=512$ code points with overlap $O=128$. An optional tokenizer-aware segmentation limits both code points and encoder tokens; dense text may require less overlap. Windows retain parent identifiers, offsets, attribution, and observation time. Stored vectors include encoder identity and source hash. Changes to weights or preprocessing require a new model identity.

Code-point limits do not imply token limits. The model card documents a 256-wordpiece truncation default \citep{minilmcard}. Our local adapter uses the embedding tokenizer, including special tokens, to check each window against its 256-token budget. It reduces window length and, when needed, overlap until the next window advances. The embedding endpoint checks the limit independently and rejects oversized inputs. Model and segmentation identities include hashes of weights, tokenizer, and preprocessing policy. A remote tokenizer remains a trusted component: the library validates its declared counts and offsets but cannot independently establish their accuracy. Queries receive the same input-limit check.

An observation supporting any superseded, archived, or expired memory is excluded from source retrieval. This conservative rule can suppress useful information when a source supports several memories. Unassociated observations can still contain obsolete claims, so eligibility alone does not establish current truth.


\subsection{Evidence collection and event dates}
A bounded evidence operation combines up to four caller-supplied query variants and deduplicates source windows by parent identifier and offset. Reciprocal rank fusion orders the combined results. The result states whether candidate lists were truncated, whether vector indexing was incomplete, and whether source changes invalidated retrieved material. Its coverage claim is limited to retrieved evidence. Several retrieved sources may describe one event, and a complete ranking list need not contain every relevant event.

Event-date resolution takes an exact source span and an explicit IANA timezone. Supported expressions include a complete calendar date, today, yesterday, a stated number of days ago, and the previous named weekday. Relative expressions use the local civil date of the source timestamp. Dates lacking a year and ambiguous week expressions remain unresolved. The result retains the source hash, original phrase, anchor and resolution rule. This validates date arithmetic, not event occurrence: a quoted or hypothetical event still needs semantic interpretation. The host owns event extraction, identity, correction handling and any final count.

\subsection{A source-checked evidence reader}
An optional reader uses three bounded model calls: question-based query planning, evidence extraction, and answer composition. One malformed response or inexact quotation may receive a single additional repair call, whose result passes the original validation. The host supplies an explicit question date and timezone and authorizes transfer of selected source excerpts. The reader fuses the existing retrieval channel with a lexical channel that normalizes simple plural endings and month names. Speaker preference and per-source limits reduce crowding by long generic responses. The serialized evidence array is bounded to 32,000 UTF-8 bytes and 80 excerpts. This is a wire-size budget, not a tokenizer guarantee or an exhaustive search claim.

An extracted fact retains an exact source quotation, source identifier, code-point offsets, source hash, observation time and a model-written interpretation. Quotations must occur within the excerpts actually supplied. Proposed date phrases must occur inside those quotations to support arithmetic; an unmatched phrase remains unresolved. Supported missing-year forms use the observation year under an explicitly reported assumption. This optional reader policy is separate from the stricter date-resolution operation, which leaves missing years unresolved.

Civil-day and completed-calendar-month differences are computed deterministically from resolved endpoints. Event counting counts the distinct event identifiers selected by the model, not source rows. The model still decides whether two accounts describe the same event, whether a visit occurred, and whether it falls inside the requested period. Exact quotes establish provenance, not semantic truth. The final answer cites ledger identifiers; its phrasing and interpretation remain statistical outputs. Cited date assumptions and inferred years produce a conditional response status, with assumptions retained separately.

All selected sources are rechecked for owner, content and lifecycle eligibility after composition. If any selected source changed or became unavailable, the operation returns an error without partial facts or an answer. This cannot retract evidence that an authorized provider already received or an answer returned before a later deletion.

\subsection{Activation and outcome utility}
For an unpinned memory with base half-life $H>0$, define
\[
 A(m,t)=2^{-d/(H(1+\log n'))},\quad
 d=\max(0,t-t_\ell),\quad n'=\max(1,n).
\]
Time reduces activation while leaving evidence confidence unchanged. Pinning or a disabled half-life sets activation to one. Explicit expiry takes precedence over pinning.

Outcome utility is separate. Starting from $U_0=1/2$, an explicitly used memory with a known binary outcome updates as
\[
 U_{N+1}=\frac{(N+2)U_N+r_{N+1}}{N+3}
        =\frac{1+\sum_{i=1}^{N+1}r_i}{N+3}.
\]
Duplicate events do not update utility again. The statistic records attributed outcomes; it does not establish that a memory caused success or update model weights.

\section{Selective classification}
The classifier estimates retention and type distributions for a candidate and its surrounding text. Training and serving use the same input format. A single exact occurrence of the candidate in its context is replaced with a marker to avoid repetition. The complete formatted input is counted with the encoder tokenizer; an oversized input is rejected before inference, without silently removing context. Its training loss is
\[
 \mathcal{L}(\theta)=
 -\sum_i\log p_\theta(r_i\mid x_i)
 -\sum_{i:k_i\text{ annotated}}\log p_\theta(k_i\mid x_i).
\]
Here $\theta$ denotes the linear heads conditional on a fixed encoder. Encoder adaptation is a separate contrastive stage using type-labeled pairs; the two cross-entropy terms are fitted afterward. The type loss excludes ignored or unresolved candidates without meaningful type labels. The experiment compares a frozen sentence encoder with linear heads against contrastive encoder adaptation followed by the same heads. Neither approach learns conflict resolution or generates summaries.

\subsection{Trained components and runtime responsibilities}
The base representation is the public \texttt{all-MiniLM-L6-v2} sentence encoder, which maps text to 384-dimensional vectors \citep{minilmcard}. Cortex does not train a foundation model. The task-specific parameters are two local multinomial logistic-regression heads and, in the adapted variant, a fine-tuned copy of the encoder. For normalized representation $h(x)$, head $j\in\{r,k\}$ returns
\[
 p_j(y\mid x)=\operatorname{softmax}\!\left((W_jh(x)+b_j)/T_j\right)_y.
\]
The frozen variant fits only $W_j,b_j$. The adapted variant first performs eighty SetFit contrastive steps using training examples with type annotations, then fits the same heads. Logistic regression uses $C=10$, seed 17, and at most 1,500 iterations. Each temperature $T_j$ is selected by development loss from $\{0.5,0.75,1,1.5,2,3,5\}$. Test labels select neither weights nor temperatures.

\begin{table}[htbp]
\centering\small
\begin{tabularx}{\linewidth}{@{}L{.23\linewidth}Y L{.26\linewidth}@{}}
\toprule
Component & Function & Learned in this study\\
\midrule
External language model & Propose annotations, consolidation operations, and answers & No weight updates\\
MiniLM encoder & Represent candidate text and context & Frozen or locally adapted\\
Retention head & Durable, temporary, ignore, or uncertain & Locally fitted parameters\\
Memory-type head & Fact, behavior, preference, relationship, project, or routine & Locally fitted parameters\\
Persistence validator & Check source identity, scope, and allowed transitions & Deterministic code\\
\bottomrule
\end{tabularx}
\caption{Training and runtime roles. The pretrained encoder and external language model are third-party components. Task-specific heads and adapted weights are experimental local artifacts, not released model checkpoints.}
\label{tab:components}
\end{table}

The retention labels encode persistence policy; the type labels describe content. Neither head extracts candidate boundaries, resolves whether two mentions refer to the same event, or decides whether a new statement contradicts a prior one. Those capabilities must be evaluated separately. The external language model still proposes consolidation operations; local classifiers have no authority over those writes.

The decision rule abstains if maximum retention probability is below $0.8$, the chosen retention is uncertain, or a retained candidate's maximum type probability is below $0.8$. These are experimental operating points. High probability on synthetic text does not establish calibration on real conversations.

Records retain source spans, grouping identifiers, label provenance, and annotation-policy version. Human annotations, designed synthetic labels, and model-generated labels remain distinguishable. Whole groups belong to training, development, or test. The loader rejects exact normalized candidate text across partitions and refuses provisional teacher labels unless explicitly enabled.

Temperature is chosen on development negative log likelihood. Test labels do not participate in fitting or temperature selection. Metrics include macro-averaged F1, multiclass Brier score, calibration error, coverage, and accepted-prediction error. Exact deduplication does not eliminate semantic leakage among templated examples.

The operational boundary remains observational: predictions are returned to the caller and never modify memory state. A host can inspect disagreement before a learned decision suppresses useful information or retains an unsupported assertion.

\section{Conditional properties}
The following results concern an abstract state machine. They assume atomic transactions, immutable observations, compliant writers, and correct owner authorization by the host. They are consequences of those assumptions, not a machine-checked proof that every implementation path satisfies them.

\begin{proposition}[Scope preservation]
If initial memories reference only observations in their own owner scope, and all source additions and update targets are checked against that scope, every reachable memory preserves the property.
\end{proposition}
\begin{proof}
Induct on committed transitions. Insertion establishes the property through source validation. Contradiction creates a checked successor and a same-owner predecessor link. Reinforcement unions the existing source set with checked sources. Other transitions remove records or change metadata without adding sources. Rejected transactions leave the preceding state unchanged.
\end{proof}

\begin{proposition}[Atomic acknowledgement]
If every consolidation writer acquires the same per-owner lock before selecting its pending snapshot, retains that lock through commit, and commits operations and acknowledgements atomically, no observation is acknowledged by two successful chunks.
\end{proposition}
\begin{proof}
The lock orders successful consolidations for an owner. At the first committed acknowledgement, the observation ceases to be pending and later snapshots exclude it. A failed transaction rolls back both operations and acknowledgements. No second successful chunk can therefore select that observation as pending.
\end{proof}
This result permits several facts to cite one observation and permits repeated inference after failure.

\begin{proposition}[Exact-replay exclusion]
Suppose imports and deletion share an owner-scoped barrier. Deletion atomically writes a content tombstone and removes the source, and subsequent imports check tombstones. Once deletion commits, importing identical source text for that owner cannot restore it.
\end{proposition}
\begin{proof}
The barrier orders an import before or after deletion. An import ordered before deletion is part of the source state examined and removed. An import ordered after deletion sees the committed tombstone and skips insertion. Identical source text cannot survive either ordering.
\end{proof}
Fingerprints are assumed collision resistant. The property excludes paraphrases, copies already returned, provider retention, and separately trained weights. In-flight reads are governed by their final eligibility snapshot; deletion cannot retract an output already authorized by an earlier snapshot.

\begin{proposition}[Activation monotonicity]
For fixed $n'\geq1$ and $H>0$, activation is nonincreasing in elapsed time. For fixed positive elapsed time it is nondecreasing in $n'$.
\end{proposition}
\begin{proof}
For $d>0$,
\[
 \frac{\partial A}{\partial d}
 =-\frac{\log2}{H(1+\log n')}A<0,\qquad
 \frac{\partial A}{\partial n'}
 =\frac{d\log2}{Hn'(1+\log n')^2}A>0.
\]
For $d\geq0$, the exponential formula gives $0<A\leq1$ and continuity at zero. Together with the derivatives, this establishes the stated weak inequalities.
\end{proof}
An older assertion is not necessarily less likely to be true.

\begin{theorem}[Window coverage]
Let source length $L>0$, width $W$, and overlap $0\leq O<W$ be integers. Generate windows starting at $j(W-O)$, stopping at the first whose end reaches $L$. They cover the source. Every contiguous span of at most $O+1$ code points lies wholly in at least one window.
\end{theorem}
\begin{proof}
Consecutive starts differ by $S=W-O\leq W$, so windows have no gap. The first starts at zero and the last reaches $L$. For a span starting at $a$, choose the latest generated start $jS\leq a$. If it is the last start, its window contains the span because it reaches $L$. Otherwise $a<(j+1)S$, hence $a\leq jS+S-1$. For length $\ell\leq O+1$,
\[
 a+\ell\leq jS+S-1+O+1=jS+W.
\]
Thus the whole span fits.
\end{proof}
Coverage does not guarantee that embedding preserves meaning or retrieval selects the window.


\begin{proposition}[Validated token-window coverage]
Let accepted windows be intervals $[a_i,b_i)$ over a source of length $L>0$, with $a_1=0$, $b_n=L$, $a_{i+1}\leq b_i$, and strictly increasing starts and ends. Their union is $[0,L)$. If a correct tokenizer reports at most $T$ tokens for each window, and the encoder uses that tokenizer and special-token policy, no accepted window exceeds $T$ encoder tokens.
\end{proposition}
\begin{proof}
The first interval covers the source prefix ending at $b_1$. Inductively, the next interval begins at or before the covered prefix ends and extends it to $b_{i+1}$. The final prefix therefore ends at $L$. The token statement follows from equality between the reported count and the encoder's count for each interval.
\end{proof}
This result guarantees complete source coverage, not containment of every multi-character fact in a single window. Variable overlap does not inherit the fixed-overlap theorem's $O+1$ span guarantee.

\begin{proposition}[Civil-date subtraction]
Represent each valid Gregorian calendar date as midnight UTC. Dividing the difference in their Unix-second representations by $86{,}400$ gives their signed difference in civil days, independent of daylight-saving transitions in the timezone used to resolve the dates.
\end{proposition}
\begin{proof}
UTC midnights in this representation are separated by exactly $86{,}400$ Unix seconds for consecutive civil dates. Summing these equal differences between the two dates gives the claimed integer. The source timezone affects which civil dates are chosen, but not subtraction of those dates.
\end{proof}
The representation does not model leap seconds or elapsed local hours. It is used for civil-day questions only.

\begin{theorem}[Selective-risk bound]
Fix a classifier and acceptance rule before testing. Assume independent test draws from the deployment distribution with correct labels. Conditional on $m>0$ accepted examples, let $\widehat R$ be their mean zero-one error and $R$ the population error conditional on acceptance. For $\delta\in(0,1)$, conditional on that accepted count, with probability at least $1-\delta$,
\[
 R\leq\min\left\{1,\widehat R+
 \sqrt{\frac{\log(1/\delta)}{2m}}\right\}.
\]
\end{theorem}
\begin{proof}
For a fixed rule, accepted examples are independent draws from the distribution conditional on acceptance. Their error indicators lie in $[0,1]$ with mean $R$. Hoeffding's inequality gives
$\Prob(R-\widehat R\geq\epsilon\mid m)\leq e^{-2m\epsilon^2}$ \citep{hoeffding}.
Substitute $\epsilon=\sqrt{\log(1/\delta)/(2m)}$ and cap at one.
\end{proof}
Selecting the rule using these test labels invalidates this experiment. The bound cannot certify deployment risk from templated synthetic data or unreviewed teacher labels.

\begin{proposition}[Source-bound answer publication]
Assume source hashes are collision-free on the compared records. Suppose every published fact is checked against a supplied source excerpt, selected sources retain their content hashes, and all selected sources pass an owner and eligibility check immediately before publication. Then every quoted fact in a successful answer is a substring of a source that was eligible at that final check.
\end{proposition}
\begin{proof}
Quote validation establishes membership in a supplied excerpt and records its parent source hash and offsets. The final check establishes that the same source content remains eligible in the requested scope. Substring membership is preserved by content equality. If either check fails, no answer is published, so a successful output has the stated property.
\end{proof}
The result is snapshot-relative. It neither verifies the model-written interpretation nor eliminates the interval between the final database read and delivery to a caller.

\section{Experimental design}
The expanded study uses two separate data sources: authorized local conversation records and a sample from LongMemEval. The local inventory includes all discovered Codex and Claude session files, with explicit user-role messages retained after provenance and content filtering. Whole project groups are assigned before classifier fitting. Earlier diagnostic studies are reported separately in Appendix~B; their scores are not pooled with the expanded results.

For public evaluation, twelve previously examined cases form the development set. Sixty additional cases are selected by a fixed hash ordering, with ten from each question type and no overlap with the nineteen cases used in earlier work. Each case receives a separate owner scope and its full supplied conversation history. Questions and reference answers never enter the memory store. Both readers use the same evidence-size ceiling, but their model-call budgets differ. The requested model is Luna 6 at medium effort through the local API. Date interpretation uses an explicit UTC reference. Scoring checks consistency with the supplied history and reference; it does not independently verify the truth of historical assistant recommendations. Model judgments are followed by Codex inspection, with no independent human adjudicator.

Classification uses the pinned MiniLM sentence encoder and the same candidate-context format for training and serving. A frozen-encoder baseline is compared with eighty steps of contrastive adaptation followed by linear heads. Temperature is chosen on development loss. Synthetic training and development examples supplement natural examples; evaluation of transfer uses only held-out natural project groups. All natural annotations remain provisional teacher labels. Test labels do not select weights or temperatures.

The local execution environment uses PostgreSQL 17 in Docker, Go 1.26.10 and a tokenizer-checked ONNX embedding adapter. A private four-slot serving variant accelerates the larger evaluation; its vectors and segmentation were checked against the serial adapter, including concurrent repeated requests. Shared machine and provider load prevents interpreting timings as an isolated throughput benchmark. Source versions, budgets and aggregate measurements are recorded in the accompanying manifest.

\section{Results}
\subsection{Expanded natural-conversation study}
A broader inventory inspected 640 local session files. After removing machine
context, generated subagent messages, detected secret patterns and exact
normalized duplicates, it retained 2,814 explicit user messages from 119 sessions
and 35 project groups: 1,336 messages from Codex and 1,478 from Claude. These
messages contained 2,111,697 characters. All eligible messages were submitted
for provisional annotation; long messages were split explicitly rather than
truncated. Annotation failures were retried in smaller batches. One retry batch
still failed and supplied no accepted labels.

The teacher produced 3,235 candidate spans. Validation excluded 192 invalid
annotations, ten exact candidate duplicates across partitions, and one duplicate
identifier. The remaining 3,032 candidates all fit the complete encoder input
budget. They contained 1,763 temporary, 535 durable, 443 ignored and 291 uncertain
labels. These are policy annotations by a model, with no human-reviewed labels.

A fixed hash ordering assigned 21 whole project groups to training, seven to
development and seven to testing, before classifier fitting. This produced
1,629, 1,040 and 363 natural candidates respectively. Synthetic training and
development partitions were added as in the earlier study. The following table
uses only the natural test partition; 310 candidates have a type label.

\begin{table}[htbp]
\centering
\begin{tabular}{lrrrr}
\toprule
Training & Retention & Type & Accepted & Errors\\
\midrule
Synthetic only, frozen & 73/363 & 186/310 & 108/363 & 78/108\\
Mixed, frozen & 215/363 & 279/310 & 61/363 & 18/61\\
Mixed, adapted & 233/363 & 285/310 & 45/363 & 13/45\\
\bottomrule
\end{tabular}
\caption{Agreement with provisional labels on seven held-out project groups.
The joint decision rule uses the fixed 0.8 thresholds. An accepted error is an
incorrect retention label or an incorrect type for a retained candidate.}
\label{tab:expanded-classification}
\end{table}

\noindent\begin{minipage}{\linewidth}
Natural examples improved agreement relative to synthetic-only training, but
this comparison does not establish useful performance against simple
constant-label policies. The adapted model reached 64.2\% retention agreement
and 91.9\% type agreement. It accepted only 12.4\% of candidates, with 28.9\%
disagreement among accepted predictions. This is insufficient evidence for
autonomous retention decisions. The apparent zero-error result among seven
accepted predictions in the smaller study did not persist at this scale. A live Go API check reproduced offline predictions for
361 candidates with matching context, with a maximum probability difference
below $10^{-6}$ and no memory writes. Two candidates at an explicitly split
source boundary had different complete context and were excluded from this
numerical-equivalence check.
\end{minipage}

\subsection{Class imbalance and calibration}
Table~\ref{tab:class-support} reports the natural test support and the adapted model's correct predictions. The dataset reflects one owner's work conversations: 253 of 363 retention labels are temporary, and 286 of 310 typed candidates are projects. Always predicting temporary would agree on 69.7\% of retention examples; always predicting project would agree on 92.3\% of typed examples. These are retrospective constant-label accuracy checks, not trained models or independently validated policies. Both exceed the adapted model's corresponding accuracy. The adapted type head predicts project for 309 of 310 inputs and preference for one; all correct type predictions are projects.

\begin{table}[htbp]
\centering\small
\begin{tabular}{@{}llrrr@{}}
\toprule
Head & Provisional label & Test support & Correct & Recall (\%)\\
\midrule
Retention & Durable & 37 & 1 & 2.7\\
 & Temporary & 253 & 213 & 84.2\\
 & Ignore & 32 & 16 & 50.0\\
 & Uncertain & 41 & 3 & 7.3\\
\midrule
Type & Fact & 9 & 0 & 0.0\\
 & Behavior & 0 & 0 & ---\\
 & Preference & 14 & 0 & 0.0\\
 & Relationship & 0 & 0 & ---\\
 & Project & 286 & 285 & 99.7\\
 & Routine & 1 & 0 & 0.0\\
\bottomrule
\end{tabular}
\caption{Class support and adapted-model recall against provisional teacher labels. Missing test classes are unevaluated. Counts are calculated from the frozen confusion matrices, not a new inference run.}
\label{tab:class-support}
\end{table}

\begin{table}[htbp]
\centering\small
\begin{tabular}{@{}lrrrrrr@{}}
\toprule
 & \multicolumn{2}{c}{Macro-F1} & \multicolumn{2}{c}{Brier} & \multicolumn{2}{c}{ECE}\\
Model & Ret. & Type & Ret. & Type & Ret. & Type\\
\midrule
Synthetic, frozen & .133 & .153 & 1.336 & .576 & .648 & .178\\
Mixed, frozen & .339 & .158 & .547 & .169 & .072 & .052\\
Mixed, adapted & .341 & .160 & .510 & .147 & .074 & .031\\
\bottomrule
\end{tabular}
\caption{Additional metrics from the same frozen predictions. Macro-F1 includes all four retention and all six declared type labels, assigning zero to undefined class scores; two type classes have no test support. Brier is multiclass squared probability error. ECE uses ten equal-width confidence bins. Lower Brier and ECE are better, but imbalance can make either look favorable while minority classes fail.}
\label{tab:classifier-calibration}
\end{table}

The type macro-F1 of .160 and zero recall on every supported non-project type qualify the 91.9\% headline accuracy. No six-type generalization claim follows from this test. The retention head also recovers only one of 37 durable examples. Although development calibration reduces aggregate probability error relative to synthetic-only training, it does not make high-confidence retention reliable. More provisional examples alone do not remedy an ambiguous annotation policy or missing class coverage. Human review, class-specific errors, and representative held-out groups are required before assigning either head operational authority.


\subsection{Operational and answer evaluation}
The full natural corpus indexed into 7,184 source windows. An observed
whitespace-only window exposed a client-side rejection that had interrupted
indexing; accepting nonempty whitespace preserved complete coverage. The final
local indexing run completed without silently truncating an input.

Large-context consolidation also exposed invalid model-generated source
references. Validation rejected those proposals and retained completed batches.
The final context uses request-local aliases and omits historical source lists.
Aliases resolve to exact IDs before the existing transaction checks. The
operational results below distinguish successful acknowledgement from semantic
correctness of the retained memories.

The recovered full-corpus run acknowledged all 2,814 observations exactly
once, with no missing or foreign source references. A fresh run using local
aliases processed 2,600 observations before a summary-length rejection; a retry
processed the remaining 214. Across the fresh run and retry, fifteen successful
chunks acknowledged all observations, and one rejected chunk acknowledged none.
The original recovery run contained five rejected chunks. These tests establish
observed persistence behavior, not the factual accuracy of every retained memory.

The twelve-case development set was evaluated twice with the frozen reader.
Source review found 11 reference-aligned answers in the first run and 12 in the
second, including respectively three and four explicitly conditional answers.
The first run abstained on an ambiguous museum-visit date. In the second, it
provided the reference value under a stated interpretation of that date. Thus
12/12 was attained in one run; the combined result was 23/24, not evidence of
perfect repeatability.

\subsection{Held-out public evaluation}
Sixty additional LongMemEval cases were selected before implementation changes:
ten per question type, ordered by a fixed SHA256 rule, excluding all nineteen
cases used in earlier development. Full supplied histories were ingested under
separate owner scopes. Questions and reference answers were excluded from
memory ingestion. The reader implementation was frozen before the final
measurement. Earlier partial executions were retained privately and were not
used to select prompts or tune against reference answers.

Both conditions received at most eighty excerpts and 32,000 serialized UTF-8
evidence bytes. The baseline used single-query retrieval and one direct reading
call. The source-checked reader used query planning, extraction and composition,
with at most one additional format repair. This is an equal evidence-size
comparison, not an equal-call-budget comparison. Scoring combined model grading
with Codex review of references and source evidence. No independent human judge
participated. Table~\ref{tab:heldout-results} reports qualified answers separately.

\begin{table}[htbp]
\centering
\begin{tabular}{lrrrrr}
\toprule
Reader & Complete & Conditional & Partial & Incorrect & Error\\
\midrule
Direct baseline & 44 & 1 & 4 & 11 & 0 \\
Source-checked & 36 & 11 & 4 & 6 & 3 \\
\bottomrule
\end{tabular}
\caption{Sixty held-out cases per condition. A conditional answer states an
assumption and gives the matching reference value. Such an answer does not
establish that the assumption is true. Reference alignment combines the first
two columns; it is distinct from unqualified correctness.}
\label{tab:heldout-results}
\end{table}

\begin{table}[htbp]
\centering
\begin{tabular}{lrr}
\toprule
Question type & Direct baseline & Source-checked\\
\midrule
Knowledge updates & 8/10 & 9/10 \\
Multi-session reasoning & 6/10 & 7/10 \\
Assistant recall & 10/10 & 8/10 \\
Personalized responses & 3/10 & 6/10 \\
User facts & 10/10 & 10/10 \\
Temporal reasoning & 8/10 & 7/10 \\
\bottomrule
\end{tabular}
\caption{Reference-aligned answers, including conditional answers, out of ten
cases per type. Partial answers and errors are excluded.}
\end{table}

The source-checked reader returned 47/60 reference-aligned answers (78.3\%),
compared with 45/60 (75.0\%) for the baseline. Complete, unconditional answers
accounted for 36/60 (60.0\%) and 44/60 (73.3\%), respectively. The alignment
measure includes the eleven and one conditional answers shown separately in
Table~\ref{tab:heldout-results}. The two-case difference is too small to establish a
reliable advantage. Personalization improved, while several assistant-recall
and temporal cases deteriorated. The reader produced more conditional answers
and three validation errors. Two errors rejected inexact quotations; one
rejected an unresolved arithmetic endpoint. Counts of distinct events and
retrieval of every required event remain weaknesses.

The final source review also found a timezone presentation inconsistency: date
arithmetic used the configured timezone, while a displayed source timestamp
could use the database driver's timezone. The release normalizes the model's
anchors to the requested timezone. It also permits one inexact quotation to use
the existing repair allowance while preserving exact-match validation. A single
literal today or yesterday phrase can supply a conditional date candidate when
the extractor omitted an anchor; this does not establish the event date. These
post-measurement fixes are covered by regression checks and targeted replay;
they are not credited with changing the frozen sixty-case score. A further
release-model pass matched eleven of the twelve development references,
missing a location detail. The earlier 12/12 result was therefore not stable
across repeated model calls.

The larger natural-label study and the public study address different tasks.
Improved agreement with provisional retention labels does not establish improved
answer quality, and reference agreement does not certify truth outside the
supplied conversation. The learned classifiers remain advisory.

These conditions measure retrieval and reading with access to original source
messages. They do not isolate the fidelity of consolidated memories or show
that compression preserves every fact needed for future questions. Equal
evidence bytes also do not imply equal model tokens, latency, or inference
cost. A consolidated-memory-only condition and matched reader budgets are
needed to attribute any gain to the memory representation itself.



\section{Discussion}
The diagnostics distinguish separate constraints. Source access can recover information omitted by summaries, but a narrow budget can still exclude a required event. Short excerpts may improve budget utilization while separating related context across windows. Classification narrows a model's decision task, but performance on designed examples does not establish transfer to natural conversations.

Observation time, event time, and validity time have different meanings. Explicit date resolution addresses arithmetic and preserves the source anchor. It does not extract every event or identify repeated descriptions of one event. Multi-session counting still requires semantic event identity and correction handling. The present classifier performs none of those operations.

Deletion exposes a boundary between external memory and learned parameters. Persistence can remove source records and derived memories. Once private text contributes to weight updates, source deletion does not itself remove that information from weights. Training artifacts remain private, and weight-training experiments use fabricated examples and, separately, owner-authorized provisional annotations. Production training needs separate consent, minimization, model-retention, and retraining policies.

Exact vector scans and explicit capacity bounds simplify small-scale operation. They do not establish suitability for large corpora. Authentication, retention policy, scheduling, and provider access remain host responsibilities. Learned mutation policies would have to preserve these boundaries.

\section{Platform context and limits of comparison}
\label{sec:platforms}
Published platform results provide reference points for the next evaluation, not a ranking against this study. Table~\ref{tab:platforms} records vendor-reported LongMemEval headline scores inspected on 22 October 2026 (UTC). We did not rerun these systems. Their source scope, readers, inference budgets, and grading differ from the sixty-case Cortex experiment.

\begin{table}[htbp]
\centering\small
\begin{tabularx}{\linewidth}{@{}L{.16\linewidth}rY@{}}
\toprule
Platform & Reported (\%) & Reported scope and protocol qualification\\
\midrule
Mem0 & 94.4 & Managed platform, top-200 retrieval; mean 6,787 context tokens. Proprietary optimizations are excluded from the open-source SDK \citep{mem0eval}.\\
Zep & 90.2 & 451/500 questions; multiple retrieval scopes and a cross-encoder reranker. Reader and judge are GPT-5.4; reader uses medium reasoning \citep{zepeval}.\\
Hindsight & 94.6 & LongMemEval-S headline on the project overview. Its recall and agentic reflection paths have different call budgets; the headline alone does not identify a matched configuration \citep{hindsightdocs}.\\
\bottomrule
\end{tabularx}
\caption{External, vendor-reported results, unverified by this study. Cortex's 78.3\% reference alignment includes conditional answers on a different sixty-case sample and is deliberately excluded from this table. No percentage-point gap or superiority claim can be inferred.}
\label{tab:platforms}
\end{table}

Mem0 describes additive temporal metadata and time-sensitive retrieval ranking \citep{mem0temporal}. Zep's Graphiti represents changing entity relationships in a temporal knowledge graph \citep{graphitidocs}. Hindsight describes semantic, keyword, graph, and temporal retrieval, followed by reranking, as well as a separate reflection operation \citep{hindsightdocs}. These are concrete design comparisons. Their benefit for a particular workload still requires an experiment.

Cortex emphasizes a small Go/PostgreSQL package with explicit source support and state transitions. It has lexical and semantic retrieval, contradiction history, activation decay, and source-linked date arithmetic. It does not yet have a validated learned retention policy, a general event-identity model, or an evaluated graph-based retrieval stage. Its engineering checks establish observed behavior under tests; the current evidence does not establish competitive answer quality, lower cost, or lower latency than the platforms above.


\section{Limitations}
The proofs apply to the stated abstract machine. Regression tests give evidence of implementation conformance in selected cases, not exhaustive verification. No theorem establishes that an extracted statement is true, a summary complete, or a probability calibrated.

The private corpus belongs to one owner and cannot be publicly reproduced. Pattern checks can accept incorrect qualifications. The original seven-case public probe was used during development. The earlier twelve-case probe is now a development set. Repeated success on those cases cannot establish generalization; the separately selected sixty-case evaluation supplies the next out-of-sample measurement. Synthetic wrapper splits share topics and vocabulary. Teacher annotations await human review, including the expanded project-separated study. These constraints require exploratory reporting and continued shadow operation.

\section{Protocol for comparative evaluation}
\label{sec:protocol}
This section specifies proposed work. No experiment or platform score in it has been measured by Cortex. Benchmark readiness means that the runs are controlled and reproducible; competitive capability requires favorable results from those runs.

\subsection{Freeze the test contract and separate prior exposure}
The first public suite should be the full 500-question LongMemEval-S release, with dataset checksum, case identifiers, source-role policy, and question timestamps fixed before running \citep{longmem,longmemrepo}. The twelve-case development set stays a regression suite. The nineteen earlier cases and sixty-case study have already influenced development or review; their union of 79 cases must be flagged. Report full-suite results for comparability and the remaining 421 cases separately as the previously unexamined subset. This separation does not establish that a foundation model never saw the public benchmark in pretraining.

Pin each implementation commit, package version, model identifier, embedding model, prompt, tokenizer, inference setting, and evaluator configuration. Archive a public protocol before inspecting new answers. Any change made after a failure creates a new experimental condition. Confirm that source ingestion excludes benchmark questions, answers, and evaluator-only evidence flags. A subsequent LoCoMo evaluation should use its official data and category definitions as a distinct transfer test \citep{locomorepo}; do not merge its score with LongMemEval.

\subsection{Compare memory representations and complete products}
Run two tracks. In the controlled-retrieval track, each system supplies evidence to the same frozen reader. In the complete-product track, each platform uses its documented answering path, with every call and cost recorded. If a managed service cannot expose or fix an internal model, report that limitation and keep it out of claims about equal-model performance. A self-hosted Graphiti run is not a reproduction of the managed Zep headline; the same distinction applies to Mem0's open-source SDK and managed platform.

\begin{table}[htbp]
\centering\small
\begin{tabularx}{\linewidth}{@{}L{.29\linewidth}Y@{}}
\toprule
Condition & Question it isolates\\
\midrule
Raw-history retrieval & What does a fixed lexical/dense baseline recover without consolidation?\\
Cortex memories only & What survives the consolidation bottleneck?\\
Cortex source windows only & What does attributable raw evidence recover?\\
Cortex combined retrieval & Does combining summaries and sources improve over either alone?\\
Classifier off versus on & Does a learned retention decision improve utility or discard required evidence?\\
Pinned competitor versions & How does the same reader perform on each retriever's evidence?\\
Oracle evidence sessions & How much error remains when the reader receives the benchmark's relevant sessions?\\
\bottomrule
\end{tabularx}
\caption{Proposed controlled conditions. Use identical questions, source access, and reader settings where the condition permits them. Oracle evidence is a labeled diagnostic and is never available to ordinary ingestion or retrieval.}
\label{tab:ablation-plan}
\end{table}

Use a common tokenizer and an 8,000-token evidence ceiling for the first controlled comparison. State truncation and packing rules in advance. Compare one-reader-call conditions separately from a bounded, four-call condition; do not grant extra planning calls to only one system. Run three answer repetitions per question at fixed settings, recording seeds where supported. Repeated answers measure inference variability and are not independent new questions. Charge ingestion, consolidation, embedding, reranking, and answer calls separately; report cold and warm p50/p95 latency, failures, and cost amortized over an explicit number of queries. Report bytes and tokens separately.

\subsection{Measure correctness, evidence, and operational behavior}
Preserve the official benchmark score using a pinned evaluator, and additionally report complete, conditional, partial, unsupported, erroneous, and correctly abstained answers under a documented human-review rubric. Do not silently count a conditional answer as established fact. Two reviewers should independently assess a predeclared stratified sample and all disputed cases, with disagreement and adjudication recorded. Reviewers should be blinded to the system identity.

Measure evidence recall against available source annotations, citation precision, answer correctness by question type, and consolidated-memory compression at a fixed downstream accuracy. Attribute failures to ingestion, consolidation, retrieval, source interpretation, or answer generation when the evidence permits. Include corrected facts, repeated event mentions, event counting, changing preferences, missing dates, and timezones. Use paired confidence intervals for answer differences; resample at the shared history or source-group level where questions are correlated. Fix the primary comparison and any multiple-comparison correction before scoring. Report inconclusive differences as inconclusive.

Keep lifecycle tests separate from semantic accuracy: retry and crash recovery, cross-owner isolation, exact-source deletion and reimport, source invalidation during an answer, context overflow, and incomplete indexing. A model-quality gain cannot compensate for losing source ownership or acknowledging an uncommitted write. Publish protocol, versions, aggregate metrics, and permitted public outputs. Private transcripts, labels, and trained weights remain outside the public repository.

\subsection{Establish labels before promoting a classifier}
Write an annotation guide that distinguishes persistence policy from content type, direct user claims from quotations or requests, and observation time from event time. Collect independently human-reviewed examples across consenting owners and source groups, with all four retention labels and all six types represented. Keep a naturally distributed test set and a separate balanced diagnostic set; report both and mark unsupported classes unevaluated. Include constant-label and frozen-encoder baselines. Test groups and owners must be excluded from fitting, replay generation, threshold selection, and calibration.

Choose confidence thresholds on development data. Before giving a classifier write authority, require a predeclared upper bound on accepted-prediction error, useful coverage, and acceptable class-specific false retention and false deletion. For illustration, a 5\% error ceiling at 95\% confidence under the stated Hoeffding bound requires at least 600 error-free accepted independent test examples. Correlated examples from one conversation do not satisfy that independence assumption. This is a proposed risk target, not a certification of the current classifier. Cases outside the validated policy remain advisory or return to human review.

\subsection{Test the sleep analogy as an ablation}
Compare direct training with source-grounded replay, matched duplicate-example training, and recombined synthetic replay. Hold the number of training examples, update steps, model capacity, and selection budget fixed. Generate every replay example exclusively from training groups. Preserve links to contributing sources and identify generated examples as synthetic; never insert a recombination as an observed user event.

Measure retention and type macro-F1, selective risk and coverage, performance on old and newly introduced domains, and downstream answer quality on untouched groups. A replay gain must exceed the matched-compute control and survive correction and source-attribution tests. This experiment can support or reject the computational analogy; it cannot establish that the implementation reproduces NREM or REM physiology.

\subsection{Criteria for a comparative claim}
A reproducible report must identify every tested condition, retain failure denominators, and publish uncertainty alongside accuracy, cost, and latency. A claim of superiority needs a predeclared paired comparison whose interval supports improvement without violating lifecycle constraints. A claim of comparable quality needs a non-inferiority margin chosen before the run and a useful cost or operational benefit. Neither criterion is met by matching a vendor's headline percentage on a different dataset slice. Until these conditions are satisfied, Cortex should be described as an experimental memory engine with tested persistence behavior and an explicit route to comparative evaluation.


\section{Conclusion}
Cortex applies a sleep-inspired separation between rapid source capture and slower integration to an inspectable memory engine. Its implemented contribution is source-preserving storage, bounded retrieval, and validated transitions. The formal results are conditional software properties; they establish neither neural fidelity nor semantic truth.

The experiments expose the next requirements. Source-assisted answering has not shown a reliable advantage over its baseline, and class imbalance explains much of the apparent classifier accuracy. Credible comparison requires the frozen protocol above, independent labels, and event-identity tests. Source-grounded replay and synthetic recombination remain hypotheses: any benefit must survive matched-compute controls, unseen groups, and source-attribution checks.

\clearpage
\appendix
\section{Implementation and reproducibility}
The accompanying manifest records implementation identity, environment, encoder revision, and artifact hashes. These identify an experiment; they do not define the scientific claims. Public source includes classification policy, training and labeling tools, and regression tests. Private conversations, per-example predictions, provider responses, database exports, and trained artifacts are excluded.

A complete observation's source windows publish in one transaction after source eligibility is checked. Retrieval checks eligibility again after provider work. The service assumes a trusted host and binds to localhost by default. Exposed deployments must supply authentication and owner authorization.

Two offsetting grading discrepancies left an earlier public total unchanged: an unsupported abstention passed for a preference question, while a response containing the correct degree subject failed because it qualified whether graduation was established. The case-level correction shows why aggregate agreement alone is inadequate.

\section{Earlier diagnostic studies}
The following pilot studies establish the earlier failure analysis and use smaller, separately described samples. Their totals are not pooled with the expanded study.
\subsection{Synthetic classification}
\begin{table}[htbp]
\centering
\begin{tabular}{lrrrr}
\toprule
Encoder & Retention F1 & Type F1 & Coverage & Error\\
\midrule
Frozen & 0.731 & 1.000 & 0.926 & 0.050\\
Contrastively adapted & 0.943 & 1.000 & 0.944 & 0.000\\
\bottomrule
\end{tabular}
\caption{Synthetic test results. Retention uses 108 examples; type uses the 84 with type annotations. Coverage and error concern the retention head at a probability threshold of 0.8, not the joint API abstention rule.}
\label{tab:classification}
\end{table}

The frozen model misclassified all 12 ignored examples in the held-out wrapper family as durable. Contrastive adaptation reduced this count to four. The adapted model abstained on those errors at the reported retention threshold. Zero observed accepted errors in this synthetic sample is not a deployment guarantee: the sampling assumptions of the selective-risk theorem are not established.

Retention Brier score was 0.148 for the frozen model and 0.032 for the adapted model. Type F1 was 1.000 for both models, but the vocabulary and facts were shared across wrapper families. These measurements demonstrate successful fitting of a designed task, not reliable classification of arbitrary conversations.

\subsection{Transfer to provisional session annotations}
The private review queue contains 84 teacher-proposed spans from the 51 authorized messages. No label in this queue has been independently reviewed by a human. The teacher labeled 66 spans temporary, nine uncertain, five ignored, and four durable. Thus the provisional distribution differs sharply from the synthetic test partition, which contains 72 durable and 12 examples of each other retention label.

\begin{table}[htbp]
\centering
\begin{tabular}{lrrr}
\toprule
Encoder & Retention agreement & Type agreement & Joint abstentions\\
\midrule
Frozen & 9/84 & 47/74 & 51/84\\
Contrastively adapted & 10/84 & 39/74 & 61/84\\
\bottomrule
\end{tabular}
\caption{Agreement with provisional teacher annotations, including predictions on abstained candidates. Type agreement uses the 74 annotated types. These are agreement measurements, not human-validated accuracy.}
\label{tab:transfer}
\end{table}

Both classifiers ran through the Go shadow API. The improvement on synthetic retention did not transfer to this diagnostic. The low agreement, distribution shift, and uncertain reference labels provide no basis for promoting either classifier to memory control. A reviewed natural-conversation corpus is required before training or evaluating an operational replacement. One of the 84 formatted classification inputs also exceeded the encoder context limit. This finding motivated exact budget validation in the follow-up study.

\subsection{Retrieval under a fixed evidence budget}
\begin{table}[htbp]
\centering
\begin{tabular}{lrrr}
\toprule
Evidence & Local checks & Public checks & Public complete\\
\midrule
Whole observations & 19/19 & 5/7 & 4/7\\
Source windows & 19/19 & 5/7 & 4/7\\
Memories and windows & 19/19 & 5/7 & 5/7\\
\bottomrule
\end{tabular}
\caption{One reader run per condition. Public checks are automatic; complete answers are judged by a subsequent Codex case inspection, not independent human review.}
\label{tab:retrieval}
\end{table}

All conditions passed the 19 local expected-answer checks. Each passed five of seven automatic public checks and failed multi-session counting and elapsed-time reasoning. Inspection marked two passing recommendation answers as partial: they named the preferred editor but did not explicitly address the requested advanced-learning focus. The combined condition did address that focus. This post hoc completeness check illustrates why a keyword pass is insufficient evidence of full answer quality.

The window representation preserved the observed local check count under the fixed budget. These measurements do not establish a general improvement over whole-source retrieval. The temporal responses abstained rather than supplying the earlier unsupported zero-day interval, but they still did not answer the question. Neither smaller excerpts nor the revised reader instruction solved the missing-evidence problem.

The frozen-memory diagnostic preceding this comparison passed 14 of the same 19 local checks using summaries alone. That result motivated source retrieval; it does not isolate a causal effect of the present window design. Replaying all 26 window queries against the final implementation produced identical retrieval outputs.

A tokenizer audit of the 3,563 retained observations found 1,165 full observations longer than the encoder's 256-wordpiece limit. The window representation contained 10,145 excerpts; 20 still exceeded that limit, with a maximum of 427 wordpieces. Thus character windows substantially reduced input-length violations in this sample but did not eliminate truncation. Token-aware segmentation remains necessary for a strict encoder-context guarantee.


\subsection{Grading audit and implementation checks}
Case inspection changed two of 21 prior public judge decisions. A memory-only answer containing the correct degree subject was accepted, while an unsupported abstention on a personalized recommendation was rejected. These changes left the aggregate memory score at 3/7 and the source and combined scores at 5/7. The audit was performed by Codex against stored questions, reference answers, and reader outputs; it is not independent human adjudication.

Regression checks cover source scope, invalid distributions, exact offsets, deletion during inference, all-or-nothing window publication, and explicit indexing. The database race suite, static checks, build, and source vulnerability scan passed. These checks support specific implementation contracts; they do not measure semantic completeness or general answer quality.

\subsection{Project-separated follow-up classification}
The follow-up sampled 63 additional explicit user-role messages from ten local project groups (25 messages from Codex and 38 from Claude). The teacher proposed 90 spans. Three retained candidates lacked a type and failed the annotation contract; they were excluded without inventing a label. The remaining 87 spans contained 67 temporary, eight ignored, eight uncertain and four durable labels. All remained provisional, with no human-reviewed examples.

A fixed hash order assigned six whole projects to training, two to development and two to testing before model fitting. This produced 43, 28 and 16 natural candidates respectively. The training and development sets combined these candidates with the corresponding synthetic partitions. The test below contains only the 16 natural candidates from held-out projects. A synthetic-only frozen encoder supplies a matched baseline using the same input format. The adapted comparison uses 80 contrastive steps. The development partition selects temperature, while the joint acceptance rule remains fixed at 0.8.

\begin{table}[htbp]
\centering
\begin{tabular}{lrrrr}
\toprule
Training & Retention & Type & Accepted & Errors\\
\midrule
Synthetic only, frozen & 2/16 & 7/14 & 3/16 & 2/3\\
Mixed, frozen & 11/16 & 13/14 & 6/16 & 1/6\\
Mixed, adapted & 10/16 & 13/14 & 7/16 & 0/7\\
\bottomrule
\end{tabular}
\caption{Agreement with provisional annotations on two held-out projects. The first two measures include abstained predictions. Errors among accepted predictions count a wrong retention label or, for retained candidates, a wrong type. These labels are not human-validated truth.}
\label{tab:natural-followup}
\end{table}

Natural examples improved agreement in this small comparison, while the adapted model accepted fewer than half the test candidates. Zero disagreements among seven accepted predictions cannot certify operational risk. The corpus is dominated by project requests and does not establish performance on other conversational distributions. All models remain advisory. A live service check classified the 16 candidates through the Go API and created no memories. A separate oversized-input check was rejected before inference.

\subsection{Frozen retrieval follow-up}
The later public probe selects two cases per question type using a fixed SHA256 ordering, excluding the original seven development cases. The twelve supplied histories contain 5,832 turns, of which 5,820 have nonempty text and are indexed. Questions and reference answers were excluded from ingestion; the reader receives only the question and retrieved evidence. Unlike the initial diagnostic, this experiment indexes source evidence directly without a consolidation pass.

Both retrieval conditions use the exact-tokenizer adapter and a 4,000-token serialized-evidence budget measured with \texttt{cl100k\_base}. The single-query condition retrieves up to 40 windows. The multi-query condition merges up to four query variants, generated from the question alone by Luna 6 at medium effort, into the same 40-window limit. The reader, also Luna 6 at medium effort, is asked to preserve attribution, apply explicit corrections and enumerate distinct events before counting. These design changes and the larger budget prevent a direct score comparison with the earlier 1,200-token experiment.

\begin{table}[htbp]
\centering
\begin{tabular}{lrrr}
\toprule
Retrieval & Complete & Partial & Incorrect\\
\midrule
Single query & 9 & 0 & 3\\
Multiple queries & 8 & 1 & 3\\
\bottomrule
\end{tabular}
\caption{Frozen public follow-up, twelve cases per condition. Codex inspected answers against the reference answers; no independent human adjudication was performed. Partial answers are not included in the complete total.}
\label{tab:retrieval-followup}
\end{table}
The single-query condition produced 9 complete answers out of 12. Multiple queries produced 8 complete answers, one inconsistent partial answer and three incorrect answers. This small probe provides no evidence of an aggregate improvement from query expansion. One expanded retrieval answered a studio-location question missed by single-query retrieval; another answer incorrectly withheld confirmation of a location despite retrieving its explicit source. A temporal answer stated both four and five months and was not counted as complete. Both conditions abstained on an elapsed-duration question even though dated departure and return evidence was present. These cases distinguish evidence availability from successful use by the reader.

Advice questions were assessed for meaningful personalization to prior context; illustrative suggestions in their reference answers were not treated as mandatory wording. The assessment remains subjective and was performed by Codex, not by an independent human. The small sample, stochastic query generation and reader responses do not establish statistical superiority of either condition.

\subsection{Input and implementation checks}
The final input audit covered 11,400 stored observations and 32,260 tokenizer-aware windows. Every source was completely covered; no accepted window exceeded the 256-token limit. The maximum observed count was 254. The same sources produced 21 over-budget windows under the fixed 512-code-point scheme. This checks preservation and input limits, not semantic retrieval quality.

All 19 prior private recall pattern checks passed with the strict adapter. These are repeated development diagnostics, not new independent questions. The final source binary reproduced all 24 public-condition evidence contexts exactly, separating documentation changes from the measured retrieval implementation. PostgreSQL race tests, Python contract tests, static checks, builds, vulnerability checks, and full-history secret scanning passed. Live checks confirmed strict rejection of an oversized classifier input, source-bound timezone resolution, and no memory writes from classification.



\clearpage
\bibliographystyle{plainnat}
\begingroup
\footnotesize
\raggedright
\setlength{\bibsep}{1pt}
\bibliography{references}
\endgroup
\end{document}
